% =============================================================================
% APPENDIX A: DATA AND HYPOTHESES
% Last updated: 2026-03-06
% =============================================================================

\section{Data and Hypotheses}

% =============================================
% A.1 MODEL ARCHITECTURE
% =============================================

\subsection{Model Architecture}
\label{app:model_arch}

The analysis is built from a structured input system containing 17 interconnected
components that form the complete hypothesis chain. The architecture is designed for
transparency and reproducibility: every numerical assumption is stored in a single
documented location, and modifying any parameter triggers automatic recalculation of
all downstream outputs.

\begin{table}[H]
\centering
\footnotesize
\caption{Model input architecture --- 17 components}
\label{tab:app_sheets}
\renewcommand{\arraystretch}{1.2}
\begin{tabular}{@{}p{3.8cm}p{2cm}p{7.5cm}@{}}
\toprule
\textbf{Component} & \textbf{Category} & \textbf{Content} \\
\midrule
Scenario definitions & Documentation & Model version history (V5--V7), scenario structure, changelog \\
Data flow architecture & Documentation & Dependency map across all input and output components \\
Global parameters & Hypotheses & 29 modifiable parameters: asset valuation, operating ratios, growth, depreciation \\
Carbon decoupling ratios & Hypotheses & Carbon/revenue trajectories by region $\times$ scenario (2025--2050) \\
Source traceability & Hypotheses & Literature references and confidence ratings for each parameter \\
Country financial data & Parameters & 120 countries: grid carbon intensity (5 horizons), tax, COGS, OPEX, WACC \\
Facility database & Core data & 13{,}205 geocoded facilities $\times$ 64 attributes \\
Efficiency matrices & Parameters & PPA factors and PUE by region $\times$ facility tier, with dynamic trajectories \\
Regional carbon projections & Trajectories & MW-weighted regional carbon intensity to 2050 \\
Aggregate results & Outputs & ClimVaR by scenario, region, channel, adaptation pathway \\
Efficiency sensitivity & Outputs & Sensitivity of results to PPA and PUE parameter variation \\
Regional decomposition & Outputs & Portfolio, ClimVaR$\times$region, facility type, channels, risk tiers \\
IEA reference mapping & Calibration & IEA source citations with section-level traceability \\
AI demand projections & Calibration & IEA electricity demand under three AI growth scenarios \\
Grid supply mix & Calibration & Regional generation mix projections to 2050 \\
Emissions cross-check & Calibration & Modeled emissions validated against IEA benchmarks \\
Extended risk parameters & Calibration & Grid delay, electricity price multipliers, supply chain data \\
\bottomrule
\end{tabular}
\end{table}

The model has undergone four documented version increments. Version~5 introduced dynamic
IEA trajectories replacing static parameters. Version~6 added regionalized asset values,
cooling ratios, and grid congestion modeling. Version~7 computed facility-level ClimVaR
across all scenarios, producing 24 output columns per facility (4 scenario combinations
$\times$ 6 channel metrics). Version~7.5 integrated three AI growth scenarios with binary
inclusion flags, enabling portfolio-level analysis under LTG, SAI, and AWB assumptions
from a single database.


% =============================================
% A.2 FACILITY DATABASE
% =============================================

\subsection{Facility Database}
\label{app:database}

\subsubsection{Database Dimensions}

The core facility database contains 13{,}205 rows and 64 columns. Each row represents
a single data center facility; each column carries one of six categories of information:

\begin{itemize}
\item \textbf{Location} (3 columns): geocoded latitude, longitude, and ISO country code
(120 countries represented).
\item \textbf{Identity and capacity} (8 columns): unique identifier, IT capacity (MW),
facility name, owner, provider type, sub-region, rack count, source database.
\item \textbf{Financial structure} (14 columns): annual revenue, asset value, COGS ratio,
OPEX ratio, tax rate, WACC, yearly growth, terminal growth, depreciation ratio, CAPEX
ratio, insurance premium, electricity cooling ratio, outdoor salary ratio, purchase ratio.
\item \textbf{Efficiency and emissions} (6 columns): Scope~1 emissions (tCO$_2$/yr),
Scope~2 raw and PPA-adjusted, PUE, PPA factor, and terminal value year.
\item \textbf{ClimVaR outputs} (24 columns): facility-level ClimVaR percentage and five
channel decompositions ($\gamma$, CC$^2$, $\beta$, $\delta$, $\xi$), under 4 scenario
combinations (BAU/NZT $\times$ MEAN/Q95).
\item \textbf{Scenario flags} (3 columns): binary inclusion flags for LTG, SAI, AWB
growth scenarios.
\end{itemize}


\subsubsection{Installed Base: Sources and Coverage}

The installed base comprises 8{,}572 facilities totaling 90{,}818~MW (90.8~GW),
assembled from four complementary source types:

\begin{enumerate}
\item \textbf{Commercial databases}: Data Center Map, Cloudscene, and proprietary
Schneider Electric client records provide the core inventory with location, operator,
and capacity attributes. Cross-referencing eliminates duplicates.
\item \textbf{Regulatory filings}: planning permissions, environmental impact assessments,
and grid connection applications in the US, EU, China, and Japan provide independent
capacity verification for facilities above 10~MW.
\item \textbf{Satellite imagery validation}: building footprint analysis provides
independent size estimates, calibrated against known facilities to derive MW capacity
from floor area.
\item \textbf{Operator disclosures}: annual and sustainability reports from publicly
listed operators (Equinix, Digital Realty, CyrusOne, NTT, Chindata, GDS) provide
facility-level data for approximately 35\% of total installed GW.
\end{enumerate}

\paragraph{Capacity distribution.} Facility IT capacities range from 15~kW (enterprise
server rooms) to 300~MW (hyperscale campuses), with a median of 10.6~MW. The distribution
is heavily right-skewed: 73.4\% of facilities exceed 1~MW, but only 4.7\% exceed 50~MW
and 3.0\% exceed 100~MW. The top decile by capacity accounts for over 60\% of total GW.

\begin{table}[H]
\centering
\small
\caption{Installed base capacity distribution}
\label{tab:app_mw_dist}
\begin{tabular}{@{}lrrr@{}}
\toprule
\textbf{Capacity band} & \textbf{Facilities} & \textbf{\% of base} & \textbf{Cumulative GW} \\
\midrule
$<$ 1 MW & 2{,}283 & 26.6\% & 0.9 \\
1--5 MW & 3{,}221 & 37.6\% & 9.5 \\
5--10 MW & 906 & 10.6\% & 6.6 \\
10--50 MW & 1{,}755 & 20.5\% & 36.5 \\
50--100 MW & 152 & 1.8\% & 10.7 \\
$>$ 100 MW & 255 & 3.0\% & 26.6 \\
\midrule
\textbf{Total} & \textbf{8{,}572} & \textbf{100\%} & \textbf{90.8} \\
\bottomrule
\end{tabular}
\end{table}

\paragraph{Provider types.} The database classifies facilities into five provider
categories reflecting operational characteristics and financial profiles.

\begin{table}[H]
\centering
\small
\caption{Installed base by provider type}
\label{tab:app_provider}
\begin{tabular}{@{}lrrl@{}}
\toprule
\textbf{Provider type} & \textbf{Facilities} & \textbf{\% of base} & \textbf{Typical capacity} \\
\midrule
Retail colocation & 5{,}662 & 66.1\% & 0.015--5~MW \\
Wholesale / large colocation & 1{,}560 & 18.2\% & 5--50~MW \\
Telecommunications & 1{,}068 & 12.5\% & 0.5--20~MW \\
Hyperscale (cloud) & 149 & 1.7\% & 50--300~MW \\
Other (hosting, reseller, cloud) & 133 & 1.6\% & Variable \\
\bottomrule
\end{tabular}
\end{table}

\paragraph{Country coverage.} Table~\ref{tab:app_country_coverage} presents the 20
largest data center markets by installed capacity. Together they account for 93\% of
global installed MW. The full database covers 120 countries.

\begin{table}[H]
\centering
\small
\caption{Top 20 data center markets by installed capacity}
\label{tab:app_country_coverage}
\begin{tabular}{@{}lrrr|lrrr@{}}
\toprule
\textbf{Country} & \textbf{DCs} & \textbf{MW} & \textbf{\%} & \textbf{Country} & \textbf{DCs} & \textbf{MW} & \textbf{\%} \\
\midrule
United States & 2{,}383 & 30{,}329 & 33.4 & Canada & 272 & 1{,}327 & 1.5 \\
China & 1{,}284 & 25{,}076 & 27.6 & Russia & 255 & 1{,}054 & 1.2 \\
United Kingdom & 374 & 3{,}596 & 4.0 & South Korea & 94 & 1{,}023 & 1.1 \\
Japan & 447 & 2{,}940 & 3.2 & Norway & 27 & 896 & 1.0 \\
India & 206 & 2{,}603 & 2.9 & Hong Kong & 109 & 809 & 0.9 \\
Germany & 291 & 2{,}128 & 2.3 & UAE & 49 & 775 & 0.9 \\
France & 186 & 1{,}682 & 1.9 & Israel & 61 & 739 & 0.8 \\
Australia & 247 & 1{,}624 & 1.8 & Italy & 151 & 716 & 0.8 \\
Singapore & 115 & 1{,}547 & 1.7 & Brazil & 116 & 683 & 0.8 \\
Ireland & 63 & 1{,}542 & 1.7 & Netherlands & 207 & 1{,}434 & 1.6 \\
\bottomrule
\end{tabular}
\end{table}

\paragraph{Aggregate emissions.} The installed base generates approximately 1.0~MtCO$_2$/yr
in Scope~1 emissions (diesel backup) and 482.7~MtCO$_2$/yr in raw Scope~2 emissions
(purchased electricity). After PPA adjustment, Scope~2 decreases to 400.0~MtCO$_2$/yr
(17\% reduction). Total annual revenue is \$147~billion; total book asset value is
\$750~billion.


\subsubsection{Prospective AI-Era Infrastructure}

Prospective facilities are drawn from four source databases, ordered by decreasing certainty.

\begin{table}[H]
\centering
\small
\caption{Prospective facility source databases}
\label{tab:app_source_databases}
\begin{tabular}{@{}p{4cm}rrrl@{}}
\toprule
\textbf{Source} & \textbf{Facilities} & \textbf{GW} & \textbf{Mean MW} & \textbf{Certainty} \\
\midrule
Epoch AI GPU Clusters & 62 & 7.8 & 125 & High \\
Tier~2 Announced & 46 & 4.6 & 100 & Medium \\
Probabilistic LTG & 2{,}316 & 231.3 & 100 & Low \\
Probabilistic SAI & 2{,}195 & 219.1 & 100 & Low \\
\bottomrule
\end{tabular}
\end{table}

\small\textit{The total unique database contains 4{,}633 prospective facilities (462.7~GW).
Facilities are assigned to growth scenarios via binary inclusion flags. The effective
count per scenario is: LTG = 615 (62.9~GW), SAI = 1{,}195 (121.1~GW), AWB = 2{,}939
(295.6~GW). Results in Parts~IV--VI are reported by scenario, not for the raw total.}
\normalsize

\paragraph{Tier~1 confirmed projects (sample).} The Epoch AI GPU Cluster database
includes named, publicly documented projects: OpenAI Stargate (Abilene TX, phased to
500~MW), Microsoft AI campuses (UK, Sweden, Spain), Meta superclusters (multiple US
locations), xAI Memphis (500~MW announced), NVIDIA Blackwell clusters (multiple sites),
AWS sovereign deployments (Saudi Arabia, Malaysia, Thailand, New Zealand, Chile), and
Google APAC expansion projects. These represent the highest-certainty prospective
capacity---most are under construction or recently commissioned.

\paragraph{Tier~2 announced projects (sample).} Includes sovereign AI campuses
(Sesterce Valence FR 250~MW, DataVolt NEOM SA 1.5~GW phased), large colocation expansions
with planning permissions (Pure DC UK, Vantage \$2B deployment, Brookfield \$10B plan),
and national programs (SK 3~GW plan South Korea, Scala 4.75~GW Brazil).


\subsubsection{Scenario Composition}

\begin{table}[H]
\centering
\small
\caption{Portfolio composition by AI growth scenario}
\label{tab:app_scenario_composition}
\begin{tabular}{@{}lrrrr@{}}
\toprule
\textbf{Scenario} & \textbf{Total DCs} & \textbf{Total GW} & \textbf{AI added} & \textbf{AI GW} \\
\midrule
Installed only & 8{,}572 & 90.8 & --- & --- \\
+ Limits to Growth & 9{,}187 & 153.7 & 615 & 62.9 \\
+ Sustainable AI & 9{,}767 & 211.9 & 1{,}195 & 121.1 \\
+ Abundance Without Boundaries & 11{,}511 & 386.4 & 2{,}939 & 295.6 \\
\bottomrule
\end{tabular}
\end{table}


\subsubsection{Regional Distribution}

\begin{table}[H]
\centering
\small
\caption{Facility coverage by region}
\label{tab:app_regional}
\begin{tabular}{@{}lrr|rr@{}}
\toprule
& \multicolumn{2}{c|}{\textbf{Installed base}} & \multicolumn{2}{c}{\textbf{Total (AWB)}} \\
\textbf{Region} & \textbf{DCs} & \textbf{GW} & \textbf{DCs} & \textbf{GW} \\
\midrule
United States & 2{,}383 & 30.3 & 4{,}071 & 199.1 \\
China & 1{,}284 & 25.1 & 2{,}267 & 123.4 \\
Europe & 2{,}019 & 15.4 & 2{,}574 & 70.9 \\
Asia-Pacific & 1{,}593 & 11.6 & 2{,}495 & 101.2 \\
Middle East & 255 & 2.6 & 370 & 14.1 \\
Rest of World & 1{,}038 & 5.9 & 1{,}428 & 44.9 \\
\midrule
\textbf{Total} & \textbf{8{,}572} & \textbf{90.8} & \textbf{11{,}511} & \textbf{386.4} \\
\bottomrule
\end{tabular}
\end{table}

\small\textit{Asia-Pacific excludes China; includes Japan, Australia, Singapore, South Korea.
Rest of World includes India, Canada, Latin America, Africa, Russia.}
\normalsize


\subsubsection{Geocoding and Climate Linkage}

All facilities are geocoded to coordinates enabling linkage to downscaled CMIP6 climate
projections at $0.1° \times 0.1°$ resolution ($\sim$10~km). For installed facilities,
coordinates are derived from address geocoding or satellite imagery. For prospective
facilities, coordinates are assigned at the announced site location (Tier~1/2) or
metropolitan area centroid (probabilistic). The climate linkage assigns four physical risk
variables to each facility for each year (2025--2050) under each RCP pathway at three
distribution quantiles (Q25, MEAN, Q95).


% =============================================
% A.3 FINANCIAL PARAMETERS
% =============================================

\subsection{Financial Parameter Calibration}
\label{app:financial}

\subsubsection{Global Parameters}

\begin{table}[H]
\centering
\footnotesize
\caption{Global model parameters (29 documented hypotheses)}
\label{tab:app_global_params}
\renewcommand{\arraystretch}{1.15}
\begin{tabular}{@{}p{4cm}rlll@{}}
\toprule
\textbf{Parameter} & \textbf{Value} & \textbf{Unit} & \textbf{Source} & \textbf{Confidence} \\
\midrule
Asset value / MW (base) & 9 & \$M/MW & Schneider, CBRE 2024 & High \\
PUE (global default) & 1.40 & --- & Uptime Institute 2024 & High \\
Backup ratio (diesel) & 0.50 & --- & Industry standard & High \\
Diesel consumption & 300 & L/MWh & Generator spec & High \\
Diesel emission factor & 2.68 & kgCO$_2$/L & GHG Protocol & High \\
AI revenue multiplier & 2.0$\times$ & --- & McKinsey, Thunder Said & Mixed \\
Size mult.\ (retail) & 1.0$\times$ & --- & Baseline & High \\
Size mult.\ (wholesale) & 0.87$\times$ & --- & CBRE, Houlihan Lokey & High \\
Size mult.\ (hyperscale) & 0.75$\times$ & --- & CBRE, Synergy Research & High \\
Revenue growth & 5\%/yr & --- & Industry consensus & Mixed \\
Terminal growth & 2\% & --- & Standard DCF & High \\
Elec.\ cooling / revenue & 5\% & --- & Uptime, IEA 2025 & Mixed \\
Outdoor salary / revenue & 0.4\% & --- & Expert estimate & Low \\
Purchase ratio ($p$) & 15\% & --- & $p \approx c - s$ & Mixed \\
Insurance premium (base) & 5\% & of damage & Industry standard & Mixed \\
Insurance premium (AI) & 10\% & of damage & Expert estimate & Low \\
Depreciation / AV ($d$) & 8\% & --- & Damodaran, Equinix & High \\
CAPEX / revenue (base) & 10\% & --- & CBRE mature DC & High \\
CAPEX / revenue (AI tiers) & 20\% & --- & Equinix 32\%, DLR 27\% & Mixed \\
\bottomrule
\end{tabular}
\end{table}


\subsubsection{Asset Valuation by Region and Tier}

\begin{table}[H]
\centering
\small
\caption{Asset value per MW (\$M/MW) by region and facility tier}
\label{tab:app_asset_value}
\begin{tabular}{@{}lccccl@{}}
\toprule
\textbf{Region} & \textbf{Base} & \textbf{Tier1} & \textbf{Tier2} & \textbf{Tier3} & \textbf{Rationale} \\
\midrule
US & 9 & 40 & 25 & 30 & Highest construction costs \\
Europe & 9 & 38 & 24 & 28 & $\sim$95\% of US costs \\
China & 7 & 28 & 18 & 22 & $\sim$70\% of US costs \\
Asia-Pacific & 8 & 32 & 20 & 25 & Mix SG/JP (high) + ID/TH (low) \\
Middle East & 9 & 42 & 26 & 32 & Extreme heat cooling premium \\
RoW & 8 & 30 & 20 & 24 & Mixed; LatAm $\sim$85\% of US \\
\bottomrule
\end{tabular}
\end{table}

\small\textit{Sources: IEA (2025); Schneider Electric TradeOff Tool; CBRE 2024. Tier1 upper
bound (\$42M/MW, Middle East) reflects IEA frontier AI estimates.}
\normalsize

\vspace{0.2cm}
The base asset value of \$9M/MW produces a total installed book value of \$750~billion.
Size multipliers and the AI revenue multiplier produce an effective range from \$6.8M/MW
(large legacy wholesale) to \$42M/MW (Tier~1 AI, Middle East). The DCF valuation ($V_0$)
exceeds book value by approximately $\times$1.35 for the installed base (\$1{,}012B on
\$750B), reflecting the going-concern premium of operational facilities.


\subsubsection{Observed Parameter Ranges}

Table~\ref{tab:app_ranges} reports the actual ranges computed across all 13{,}205
facilities, validating that calibration produces realistic distributions.

\begin{table}[H]
\centering
\small
\caption{Observed parameter ranges across the database}
\label{tab:app_ranges}
\begin{tabular}{@{}lrrrr@{}}
\toprule
\textbf{Parameter} & \textbf{Min} & \textbf{Median} & \textbf{Max} & \textbf{Unit} \\
\midrule
IT capacity & 0.015 & 11.2 & 300.0 & MW \\
Asset value & \$120K & \$91M & \$8.0B & \$ \\
Annual revenue & \$29K & \$18.8M & \$720M & \$/yr \\
COGS / revenue & 35\% & 46\% & 50\% & ratio \\
OPEX / revenue & 22\% & 27\% & 30\% & ratio \\
WACC & 5.5\% & 7.0\% & 15.5\% & ratio \\
Tax rate & 0\% & 25\% & 35\% & ratio \\
PUE & 1.12 & 1.40 & 1.55 & ratio \\
PPA factor & 0.49 & 0.70 & 0.90 & ratio \\
ClimVaR\% (BAU/MEAN) & 2.4\% & 8.8\% & 16.3\% & \% of $V_0$ \\
ClimVaR\% (NZT/MEAN) & 3.0\% & 9.7\% & 17.0\% & \% of $V_0$ \\
\bottomrule
\end{tabular}
\end{table}


\subsubsection{Country-Level Parameters}

The model provides 14 parameters for each of 120 countries: grid carbon intensity at five
time horizons (2024, 2030, 2035, 2040, 2050), tax rate, COGS ratio, OPEX ratio, WACC, and
regional classification. Table~\ref{tab:app_country_full} presents the 20 largest markets.

\begin{table}[H]
\centering
\footnotesize
\caption{Financial and carbon parameters --- 20 largest data center markets}
\label{tab:app_country_full}
\begin{tabular}{@{}l r rrrrr rrr@{}}
\toprule
& & \multicolumn{5}{c}{\textbf{Grid carbon intensity (gCO$_2$/kWh)}} & & & \\
\cmidrule(lr){3-7}
\textbf{Country} & \textbf{Region} & \textbf{2024} & \textbf{2030} & \textbf{2035} & \textbf{2040} & \textbf{2050} & \textbf{Tax} & \textbf{COGS} & \textbf{WACC} \\
\midrule
US & N.\ Am. & 390 & 293 & 215 & 156 & 120 & 21\% & 48\% & 6.8\% \\
CN & AP Em. & 580 & 464 & 348 & 232 & 136 & 25\% & 38\% & 9.0\% \\
DE & W.\ Eur. & 380 & 247 & 160 & 106 & 46 & 30\% & 49\% & 6.5\% \\
GB & W.\ Eur. & 220 & 143 & 93 & 62 & 26 & 25\% & 47\% & 7.2\% \\
JP & AP Dev. & 470 & 329 & 235 & 188 & 141 & 30\% & 48\% & 6.5\% \\
FR & W.\ Eur. & 55 & 36 & 23 & 15 & 7 & 25\% & 44\% & 6.8\% \\
SG & AP Dev. & 420 & 294 & 210 & 168 & 126 & 17\% & 45\% & 6.0\% \\
IN & S.\ Asia & 710 & 533 & 391 & 284 & 178 & 25\% & 35\% & 11.5\% \\
AU & AP Dev. & 620 & 496 & 372 & 260 & 136 & 30\% & 44\% & 7.2\% \\
NL & W.\ Eur. & 310 & 202 & 131 & 87 & 37 & 26\% & 47\% & 6.5\% \\
BR & Lat.\ Am. & 90 & 74 & 57 & 45 & 25 & 34\% & 40\% & 11.5\% \\
KR & AP Dev. & 450 & 315 & 225 & 180 & 90 & 24\% & 44\% & 6.5\% \\
IE & W.\ Eur. & 300 & 195 & 126 & 84 & 36 & 12\% & 47\% & 7.0\% \\
AE & M.\ East & 400 & 352 & 288 & 232 & 140 & 9\% & 36\% & 8.0\% \\
SE & W.\ Eur. & 15 & 10 & 6 & 4 & 3 & 21\% & 47\% & 6.5\% \\
NO & W.\ Eur. & 10 & 7 & 4 & 3 & 2 & 22\% & 47\% & 6.5\% \\
MY & AP Em. & 550 & 413 & 303 & 220 & 110 & 24\% & 38\% & 8.5\% \\
ID & AP Em. & 700 & 560 & 420 & 280 & 168 & 22\% & 35\% & 10.0\% \\
SA & M.\ East & 500 & 440 & 360 & 290 & 175 & 20\% & 40\% & 8.5\% \\
MX & Lat.\ Am. & 440 & 330 & 242 & 176 & 110 & 30\% & 42\% & 10.0\% \\
\bottomrule
\end{tabular}
\end{table}

\small\textit{Grid carbon intensity: IEA (2025) regional projections. Tax: OECD (2024).
COGS: Equinix, Digital Realty, CyrusOne reports. WACC: Damodaran (2024) REIT sector
with country risk premiums. The 100 remaining countries follow the same structure.}
\normalsize


% =============================================
% A.4 PHYSICAL RISK DATA
% =============================================

\newpage
\subsection{Physical Risk Data}
\label{app:physical}

\subsubsection{Climate Projection Source}

Physical risk variables are extracted from bias-corrected, statistically downscaled CMIP6
projections \citep{eyring2016} at $0.1° \times 0.1°$ resolution ($\sim$10~km at
mid-latitudes). Projections are generated at each facility's geocoded coordinates for
2025--2050 under RCP~2.6, 4.5, and 8.5, at three distribution quantiles (Q25, MEAN, Q95).

\subsubsection{Variable Definitions}

\begin{table}[H]
\centering
\small
\caption{Physical risk variables: definitions and channel mapping}
\label{tab:app_physical_vars}
\begin{tabular}{@{}p{2.2cm}p{1.5cm}p{4cm}p{5.5cm}@{}}
\toprule
\textbf{Variable} & \textbf{Channel} & \textbf{Definition} & \textbf{Climate driver} \\
\midrule
COOLING & $\gamma$ & Fractional increase in cooling energy vs.\ 2024 & Cooling degree days (CDD, base 18°C, ASHRAE) \\
BID & $\beta$ & Business interruption days/yr, fraction of year & Composite: flood, storm, drought, wildfire, heat \\
DAMAGE & $\delta$ & Annual damage probability, fraction of AV & Flood depth, wind speed, seismic, fire \\
HPL & $\xi$ & Fractional outdoor labor productivity loss & Days $>$ 35°C / 40°C wet-bulb \\
\bottomrule
\end{tabular}
\end{table}

\subsubsection{Hazard Sub-Decomposition}

For the business interruption and damage channels, the model further decomposes risk into
five hazard types: drought (d), extreme wind (ew), riverine flood (r), storm surge (s),
and tropical wind (tw). This decomposition provides the inputs for the Shapley attribution
(Section~2.2.3) and is reported at facility level in the case study outputs. Each hazard
type carries its own annual probability distribution at facility coordinates.

\subsubsection{Dataset Scale}

The complete projection dataset contains approximately 17~million observations:
up to 11{,}511 facilities (AWB) $\times$ 25 years $\times$ 2--3 RCP pathways $\times$
4 physical risk variables $\times$ 3 quantiles. Projections are computed for all facilities
regardless of scenario inclusion, enabling flexible portfolio composition at analysis stage.


% =============================================
% A.5 SUPPLY CHAIN DATA
% =============================================

\subsection{Supply Chain Data}
\label{app:supply_chain}

\subsubsection{Leontief Input-Output Tables}

The technical coefficient matrix $A$ is drawn from OECD Multi-Regional Input-Output Tables
(MRIOT, 2019 vintage): 67 countries $\times$ 45 ISIC sectors. The data center sector maps
to ISIC telecommunications, computer services, and electrical equipment manufacturing.

\subsubsection{Supply Chain Concentration}

\begin{table}[H]
\centering
\small
\caption{Critical supply chain components for data center infrastructure}
\label{tab:app_supply_chain}
\begin{tabular}{@{}p{3cm}cp{2.5cm}p{2cm}p{3.5cm}@{}}
\toprule
\textbf{Component} & \textbf{China share} & \textbf{Category} & \textbf{Substitut.} & \textbf{DC impact} \\
\midrule
GaN devices & 99\% & Semiconductors & Very low & Power conversion \\
High-purity Si (4N+) & 95\% & Raw material & Very low & Semiconductor base \\
Rare earth magnets & 90\% & Cooling & Very low & Fans, pumps, compressors \\
LFP battery cells & 99\% & Storage & Low & UPS, BESS \\
SiC substrates & 50\% & Semiconductors & Medium & Power supply \\
HBM memory & $\sim$40\% (KR) & Memory & Low & GPU bandwidth \\
Refined copper & 40\% & Cabling/cooling & Medium & Busbars, heat exchangers \\
Transformers & $\sim$35\% & Grid equipment & Medium & DC power feed \\
\bottomrule
\end{tabular}
\end{table}

\small\textit{Sources: IEA (2025) Table 4.1; IEA Critical Minerals; USGS.}
\normalsize

\subsubsection{Pass-Through Calibration}

\begin{table}[H]
\centering
\small
\caption{Leontief pass-through rates by supply chain pathway}
\label{tab:app_passthrough}
\begin{tabular}{@{}p{3.5cm}ccp{5cm}@{}}
\toprule
\textbf{Pathway} & \textbf{Rate} & \textbf{Lead time} & \textbf{Rationale} \\
\midrule
Semiconductor fab & 60--80\% & 6--12 mo & TSMC 54\%, limited substitutability \\
Power transformers & 40--60\% & 18--36 mo & Moderate concentration, long lead \\
Cooling components & 20--40\% & 3--6 mo & More substitutable \\
Server/rack assembly & 30--50\% & 3--9 mo & Component bottlenecks \\
\bottomrule
\end{tabular}
\end{table}

\small\textit{Calibrated from: \citet{barrot2016}, IEA \citep{iea2025energyai}, expert assessment.
Fixed-coefficient limitation discussed in Section~5.10.}
\normalsize


% =============================================
% A.6 OUTPUT DATA STRUCTURE
% =============================================

\subsection{Output Data Structure}
\label{app:outputs}

The model produces three output datasets, documented in Table~\ref{tab:app_outputs}.

\begin{table}[H]
\centering
\footnotesize
\caption{Model output datasets}
\label{tab:app_outputs}
\renewcommand{\arraystretch}{1.2}
\begin{tabular}{@{}p{2.8cm}p{1.3cm}p{1.3cm}p{8cm}@{}}
\toprule
\textbf{Dataset} & \textbf{Rows} & \textbf{Cols} & \textbf{Content} \\
\midrule
Aggregate results & 288 & 14 & ClimVaR by scenario (BAU/NZT) $\times$ quantile (Q25/MEAN/Q95) $\times$ region (4) $\times$ asset class (4) $\times$ adaptation (3). Columns: $V_0$, five channels, residual value, ClimVaR\%, ClimVaR. \\[0.2cm]
Annual decomposition & 288 & 305 & Same cross-section with annual resolution 2026--2050. 12 channel variables $\times$ 25 years + 3 annual summary metrics $\times$ 25 years. Enables trajectory analysis. \\[0.2cm]
Facility case studies & 72 & 592 & Three sites (Abilene, Valence, Guangdong) $\times$ 24 scenario combinations. 64 base columns + 528 time-series columns: 22 channel-hazard decompositions $\times$ 25 years. \\
\bottomrule
\end{tabular}
\end{table}

The 12 channel variables decompose ClimVaR into: direct business interruption, supply
chain interruption, physical damage (by 5 hazard types: drought, extreme wind, riverine
flood, storm surge, tropical wind), heat productivity, cooling cost increase, and carbon
costs (Scopes 1, 2, 3). The facility case study dataset further decomposes damage into the
same five hazard types, enabling the Shapley attribution presented in Section~4.3.


% =============================================
% A.7 PARAMETER CONFIDENCE
% =============================================

\subsection{Parameter Confidence and Sensitivity}
\label{app:confidence}

Of the 29 documented global parameters, 14 (48\%) are calibrated directly from published
data with high confidence, 10 (34\%) combine literature and expert judgment, and 5 (17\%)
rely primarily on expert estimates. The low-confidence parameters are minor contributors:
a $\pm$50\% perturbation of all five changes aggregate ClimVaR by $<$3\%.

\subsubsection{Validation Cross-Checks}

\paragraph{Industry benchmark.} The WEF projects climate-driven costs for data centers of
\$81~billion annually by 2035 \citep{wef2025}. Our installed base ClimVaR of \$388~billion,
annualized over 25 years at 7\% discount rate, implies approximately \$34~billion/yr---lower
than the WEF figure (which includes prospective infrastructure). Under Abundance Without
Boundaries, our annualized equivalent rises to $\sim$\$150~billion. Directional alignment
confirms that ClimVaR values are conservative.

\paragraph{Sensitivity bounds.} WACC $\pm$2~pp shifts aggregate ClimVaR by $\pm$15--20\%;
halving BI frequency reduces it by $\sim$25\%; $\pm$50\% on the five low-confidence
parameters changes aggregate ClimVaR by $<$3\%. The order of magnitude ($\sim$\$400B
installed) is robust to reasonable parameter variation. Monte Carlo sensitivity analysis
is planned for a subsequent version.